\begin{thebibliography}{99}\small\setlength{\itemsep}{2pt}\setlength{\parskip}{0pt}
\bibitem{unesco}UNESCO. Recomendación sobre la Ciencia Abierta. París (2021). \url{https://www.unesco.org/en/legal-affairs/recommendation-open-science}.

\bibitem{wenger}Wenger, E. \emph{Communities of Practice: Learning, Meaning, and Identity}. Cambridge University Press (1998). \url{https://doi.org/10.1017/CBO9780511803932}.

\bibitem{hacklearning}Schulten, C. y Chounta, I.-A. How do we learn in and from Hackathons? A systematic literature review. \emph{Education and Information Technologies} \textbf{29}, 20103--20134 (2024). \url{https://doi.org/10.1007/s10639-024-12668-1}.

\bibitem{hackmapping} Costa, G., García-Holgado, A. y Alvarez, P. Systematic mapping of the literature on hackathons. \emph{Frontiers in Education} \textbf{10}, 1701443 (publicado el 22 de enero de 2026). \url{https://doi.org/10.3389/feduc.2025.1701443}.

\bibitem{coafina}Camacho Toro, R., Briceño, Y. y Martínez, A. Hackathon Co-Afina 2023: Latin American Students Drive Innovation with Open Data. Software Sustainability Institute (7 de diciembre de 2023). \url{https://www.software.ac.uk/blog/hackathon-co-afina-2023-latin-american-students-drive-innovation-open-data}.

\bibitem{coafina2024}Briceño, Y. CoAfina returns: The Latin American hackathon focused on open data. LA-CoNGA physics (30 de mayo de 2024). \url{https://laconga.redclara.net/en/coafina-returns-the-latin-american-hackathon-focused-on-open-data/}.

\bibitem{salud}Sánchez, A., Bennett, A., Stifano, A., Peña-Rosas, J.P. y Mehta, S. Editatón Salud para todos: modelo participativo de traducción del conocimiento y movilización comunitaria en Wikipedia. \emph{Archivos Latinoamericanos de Nutrición} \textbf{76}(1), 49--56 (2026). \url{https://doi.org/10.37527/2026.76.1.006}.

\bibitem{atlas}Sánchez Pineda, A. Boosting high-energy physics education around the world with ATLAS Open Data. ATLAS/CERN (26 de julio de 2018). \url{https://atlas.cern/updates/blog/HEP-education-worldwide-with-atlas-open-data}.

\bibitem{coafinacontinuidad}Briceño, Y. Latin American students develop innovative solutions with open data at Co-Afina 2023. LA-CoNGA physics (30 de noviembre de 2023). \url{https://laconga.redclara.net/en/latin-american-students-develop-innovative-solutions-with-open-data-at-co-afina-2023/}.

\bibitem{wikihealth}Heilman, J.M., Kemmann, E., Bonert, M. y colaboradores. Wikipedia: A Key Tool for Global Public Health Promotion. \emph{Journal of Medical Internet Research} \textbf{13}(1), e14 (2011). \url{https://doi.org/10.2196/jmir.1589}.

\bibitem{openglam}OpenGLAM. Effective Open Data \& Soft: Knowledge Transfer that Boosts Computer Skills in the Americas. Anuncio de conferencia de Arturo Sánchez, 27 de abril de 2020. \url{https://openglam.at/arturo_sanchez/}.

\bibitem{hsftraining}Malik, S., Meehan, S., Lieret, K. y colaboradores. Software Training in HEP. \emph{Computing and Software for Big Science} \textbf{5}, 22 (2021). \url{https://doi.org/10.1007/s41781-021-00069-9}. Versión de autor: \url{https://arxiv.org/abs/2103.00659}.

\bibitem{cc0}Creative Commons. Public Domain: CC0 Public Domain Dedication. Documentación institucional. Consulta: 21 de octubre de 2026. \url{https://creativecommons.org/public-domain/}.
\end{thebibliography}
